\documentclass{article}
\usepackage[T1]{fontenc}
\usepackage{lmodern}
\usepackage{graphicx}
\usepackage{amsmath,amssymb}
\usepackage[a4paper,margin=2cm]{geometry}
\usepackage[absolute,overlay]{textpos}
\setlength{\TPHorizModule}{1mm}
\setlength{\TPVertModule}{1mm}
\usepackage[indonesian]{babel}
\usepackage{hyperref}
\setlength{\emergencystretch}{2em}
% unit: Halaman 8

\begin{document}

% === Halaman 8 (python/native) ===

8

1. Pilih dua bilangan prima, p dan q  (sebaiknya p  q)

2. Hitung n = pq. 

3. Hitung (n) = (p – 1)(q – 1). 

4. Pilih sebuah bilangan bulat e sebagai kunci publik, e harus relatif prima 

terhadap (n) .  

5. Hitung kunci dekripsi, d, dengan kekongruenan



 ed  1 (mod (n)) atau d  e–1 ( mod ((n) ) → 𝑑=

1+𝑘∅(𝑛)

𝑒

Hasil dari algoritma di atas:

 -  Kunci publik adalah pasangan (e, n) 

  - Kunci privat adalah pasangan (d, n)


Prosedur Pembangkitan Sepasang Kunci

\end{document}
